\begin{abstract}
The maximum flow network interdiction problem (MFNIP) asks which arcs of a capacitated network an interdictor with a limited budget should remove to minimize the maximum flow that remains. It models how the loss of critical routes reduces the capacity of a supply chain. We propose a quantum interface that applies the quantum minimum-finding algorithm of Dürr and Høyer to the cardinality-constrained MFNIP. The interface assigns one qubit to each arc, prepares a superposition of Dicke states so that the search covers only the $N$ feasible interdiction plans, and uses a phase oracle that compares the follower's maximum flow with a threshold. We prove that, with this interface, the algorithm finds an optimal plan with probability at least one half using $\mathcal{O}(\sqrt{N})$ applications of polynomial-size circuits, whereas any classical algorithm that only evaluates the flow of candidate plans needs $\Omega(N)$ evaluations. In a query-level simulation on 30 layered networks with 9 to 974{,}121 feasible plans, the mean cost of minimum finding grew as $N^{0.494}$ (95\% confidence interval $[0.460,0.523]$) and that of a classical random-order search as $N^{0.975}$, and every run found the optimum within the bound of the proof. We also report qubit and gate counts of the state preparation and of an arithmetic oracle for networks with up to 14 arcs. On the IBM Heron processor \texttt{ibm\_fez}, amplitude amplification for a three-arc network returned the optimal plan in 81.5\% of the shots, against 25\% under uniform sampling. The speedup is quadratic, asymptotic, and relative to exhaustive search over the feasible plans. On every instance we tested, a mixed-integer programming solver was faster by several orders of magnitude.
\end{abstract}

\begin{IEEEkeywords}
network interdiction, maximum flow, Grover search, quantum minimum finding, Dicke states, query complexity
\end{IEEEkeywords}
